\documentclass[10pt,a4paper]{amsart}
\usepackage{amsmath,amssymb,amsthm,graphicx}
% For dealing with references we use the comment environment
\usepackage{verbatim}
\newenvironment{reference}{\comment}{\endcomment}
%\newenvironment{reference}{}{}
\newenvironment{slogan}{\comment}{\endcomment}
\newenvironment{history}{\comment}{\endcomment}
% For commutative diagrams we use Xy-pic
\usepackage[all]{xy}

% We use 2cell for 2-commutative diagrams.
\xyoption{2cell}
\UseAllTwocells
\theoremstyle{plain}
\newtheorem{theorem}[subsection]{Theorem}
\newtheorem{proposition}[subsection]{Proposition}
\newtheorem{lemma}[subsection]{Lemma}

\theoremstyle{definition}
\newtheorem{definition}[subsection]{Definition}
\newtheorem{example}[subsection]{Example}
\newtheorem{exercise}[subsection]{Exercise}
\newtheorem{situation}[subsection]{Situation}

\theoremstyle{remark}
\newtheorem{remark}[subsection]{Remark}
\newtheorem{remarks}[subsection]{Remarks}

\numberwithin{equation}{subsection}
% Macros
%
\def\lim{\mathop{\mathrm{lim}}\nolimits}
\def\colim{\mathop{\mathrm{colim}}\nolimits}
\def\Spec{\mathop{\mathrm{Spec}}}
\def\Hom{\mathop{\mathrm{Hom}}\nolimits}
\def\Ext{\mathop{\mathrm{Ext}}\nolimits}
\def\SheafHom{\mathop{\mathcal{H}\!\mathit{om}}\nolimits}
\def\SheafExt{\mathop{\mathcal{E}\!\mathit{xt}}\nolimits}
\def\Sch{\mathit{Sch}}
\def\Mor{\mathop{\mathrm{Mor}}\nolimits}
\def\Ob{\mathop{\mathrm{Ob}}\nolimits}
\def\Sh{\mathop{\mathit{Sh}}\nolimits}
\def\NL{\mathop{N\!L}\nolimits}
\def\CH{\mathop{\mathrm{CH}}\nolimits}
\def\proetale{{pro\text{-}\acute{e}tale}}
\def\etale{{\acute{e}tale}}
\def\QCoh{\mathit{QCoh}}
\def\Ker{\mathop{\mathrm{Ker}}}
\def\Im{\mathop{\mathrm{Im}}}
\def\Coker{\mathop{\mathrm{Coker}}}
\def\Coim{\mathop{\mathrm{Coim}}}

% Boxtimes
%
\DeclareMathSymbol{\boxtimes}{\mathbin}{AMSa}{"02}

%
% Macros for moduli stacks/spaces
%
\def\QCohstack{\mathcal{QC}\!\mathit{oh}}
\def\Cohstack{\mathcal{C}\!\mathit{oh}}
\def\Spacesstack{\mathcal{S}\!\mathit{paces}}
\def\Quotfunctor{\mathrm{Quot}}
\def\Hilbfunctor{\mathrm{Hilb}}
\def\Curvesstack{\mathcal{C}\!\mathit{urves}}
\def\Polarizedstack{\mathcal{P}\!\mathit{olarized}}
\def\Complexesstack{\mathcal{C}\!\mathit{omplexes}}
% \Pic is the operator that assigns to X its picard group, usage \Pic(X)
% \Picardstack_{X/B} denotes the Picard stack of X over B
% \Picardfunctor_{X/B} denotes the Picard functor of X over B
\def\Pic{\mathop{\mathrm{Pic}}\nolimits}
\def\Picardstack{\mathcal{P}\!\mathit{ic}}
\def\Picardfunctor{\mathrm{Pic}}
\def\Deformationcategory{\mathcal{D}\!\mathit{ef}}
\usepackage{cleveref}
\crefname{originallemma}{original lemma}{original lemmas}
\Crefname{originallemma}{Original lemma}{Original lemmas}
\makeatletter
\newlabel{lemma-describe-divided-power-envelope}{{07HC}{source}{}{sourcepointer-1}{}}
\newlabel{lemma-describe-divided-power-envelope@cref}{{[originallemma][0][]07HC}{[1][0][]source}{}{}{}}
\makeatother
\crefname{originallemma}{original lemma}{original lemmas}
\Crefname{originallemma}{Original lemma}{Original lemmas}
\makeatletter
\newlabel{algebra-lemma-flat-eq}{{00HK}{source}{}{sourcepointer-2}{}}
\newlabel{algebra-lemma-flat-eq@cref}{{[originallemma][0][]00HK}{[1][0][]source}{}{}{}}
\makeatother
\title{Lemma flat extension divided power envelope in crystalline.tex}
\author{The Stacks Project Authors}
\date{stacks-project:a04446e57ec1fbc252a871afcec7752fb2807b14; GNU Free Documentation License 1.2 (as published in the Stacks Project COPYING)}
\begin{document}
\maketitle
\noindent\textit{Editorial note.} All mathematical statement and proof text is reproduced verbatim from the author's chapter crystalline.tex of the Stacks Project, fixed at commit a04446e57ec1fbc252a871afcec7752fb2807b14. The pinned preamble is supplied as its own file; the wrapper class, the xr-hyper wiring and the \textbackslash{}externaldocument list are omitted because the assembler supplies the class.\par\medskip
\noindent\textit{Editorial note.} Numbering inside this standalone excerpt is editorial, not the full Stacks book print numbering.\par\medskip
\noindent\textit{Editorial note.} The candidate is from the extractor's stacks-candidates line, not from the earlier judged initial\_corpus line; the H2 judgement and its reason below were made in this module by reading the author's own statement and complete proof in the fixed checkout. Original candidate key: stacks:crystalline.tex:lemma-flat-extension-divided-power-envelope.\par\medskip
\noindent\textit{Editorial note.} Reference labels in the proof resolve to navigation pointers to the original author source; these pointers do not reproduce proofs of the prerequisites.\par\medskip
\noindent\textit{Allowed standard prerequisites (editorial).}\begin{itemize}
\item Prerequisites named by the author inside the counted span. The counted span is reproduced verbatim, including the author's own citations. Its prerequisites are the earlier results of the same fixed Stacks Project version that the author cites there and in the chapter's lead-in.
\end{itemize}\medskip
\ifdefined\hypertarget\hypertarget{sourcepointer-1}{}\fi
\noindent\textit{Source reference.} Original lemma 07HC: sources/crystalline.tex, source lines 184--184. Navigation pointer used in the counted span; the span contains the pointer only and the cited result is not reproduced inside it. Registered as a non-load-bearing navigation stub in the established form for Stacks material. No bibliography citation is involved.\par\medskip
\ifdefined\hypertarget\hypertarget{sourcepointer-2}{}\fi
\noindent\textit{Source reference.} Original lemma 00HK: sources/algebra.tex, source lines 8995--8995. Navigation pointer used in the counted span; the span contains the pointer only and the cited result is not reproduced inside it. Registered as a non-load-bearing navigation stub in the established form for Stacks material. No bibliography citation is involved.\par\medskip
\begin{lemma}
\label{lemma-flat-extension-divided-power-envelope}
Let $(B, I, \gamma) \to (B', I', \gamma')$ be a homomorphism of
divided power rings. Let $I \subset J \subset B$ and
$I' \subset J' \subset B'$ be ideals. Assume
\begin{enumerate}
\item $B/I \to B'/I'$ is flat, and
\item $J' = JB' + I'$.
\end{enumerate}
Then the canonical map
$$
D_{B, \gamma}(J) \otimes_B B' \longrightarrow D_{B', \gamma'}(J')
$$
is an isomorphism.
\end{lemma}

\begin{proof}
Set $D = D_{B, \gamma}(J)$. Choose elements $f_t \in J$ which generate $J/I$.
Set $\mathcal{R} = \{(r_0, r_t) \in I \oplus \bigoplus\nolimits_{t \in T} B
\mid \sum r_t f_t = r_0 \text{ in }B\}$ as in the proof of
Lemma \ref{lemma-describe-divided-power-envelope}. This lemma shows that
$$
D = B\langle x_t \rangle/ K
$$
where $K$ is generated by the elements $x_t - f_t$ and
$\delta_n(\sum r_t x_t - r_0)$ for $(r_0, r_t) \in \mathcal{R}$.
Thus we see that
\begin{equation}
\label{equation-base-change}
D \otimes_B B' = B'\langle x_t \rangle/K'
\end{equation}
where $K'$ is generated by the images in $B'\langle x_t \rangle$
of the generators of $K$ listed above. Let $f'_t \in B'$ be the image
of $f_t$. By assumption (1) we see that the elements $f'_t \in J'$
generate $J'/I'$ and we see that $x_t - f'_t \in K'$. Set
$$
\mathcal{R}' =
\{(r'_0, r'_t) \in I' \oplus \bigoplus\nolimits_{t \in T} B'
\mid \sum r'_t f'_t = r'_0 \text{ in }B'\}
$$
To finish the proof we have to show that
$\delta'_n(\sum r'_t x_t - r'_0) \in K'$ for
$(r'_0, r'_t) \in \mathcal{R}'$, because then the presentation
(\ref{equation-base-change}) of $D \otimes_B B'$ is identical
to the presentation of $D_{B', \gamma'}(J')$ obtain in
Lemma \ref{lemma-describe-divided-power-envelope} from the generators $f'_t$.
Suppose that $(r'_0, r'_t) \in \mathcal{R}'$. Then
$\sum r'_t f'_t = 0$ in $B'/I'$. As $B/I \to B'/I'$ is flat by
assumption (1) we can apply the equational criterion of flatness
(Algebra, Lemma \ref{algebra-lemma-flat-eq}) to see
that there exist an $m > 0$ and
$r_{jt} \in B$ and $c_j \in B'$, $j = 1, \ldots, m$ such
that
$$
r_{j0} = \sum\nolimits_t r_{jt} f_t  \in I \text{ for } j = 1, \ldots, m
$$
and
$$
i'_t  = r'_t - \sum\nolimits_j c_j r_{jt} \in I' \text{ for all }t
$$
Note that this also implies that
$r'_0 = \sum_t i'_t f_t + \sum_j c_j r_{j0}$.
Then we have
\begin{align*}
\delta'_n(\sum\nolimits_t r'_t x_t  - r'_0)
& =
\delta'_n(
\sum\nolimits_t i'_t x_t +
\sum\nolimits_{t, j} c_j r_{jt} x_t -
\sum\nolimits_t i'_t f_t -
\sum\nolimits_j c_j r_{j0}) \\
& =
\delta'_n(
\sum\nolimits_t i'_t(x_t - f_t) +
\sum\nolimits_j c_j (\sum\nolimits_t r_{jt} x_t  - r_{j0}))
\end{align*}
Since $\delta_n(a + b) = \sum_{m = 0, \ldots, n} \delta_m(a) \delta_{n - m}(b)$
and since $\delta_m(\sum i'_t(x_t - f_t))$ is in the ideal
generated by $x_t - f_t \in K'$ for $m > 0$, it suffices to prove that
$\delta_n(\sum c_j (\sum r_{jt} x_t  - r_{j0}))$ is in $K'$.
For this we use
$$
\delta_n(\sum\nolimits_j c_j (\sum\nolimits_t r_{jt} x_t  - r_{j0}))
=
\sum c_1^{n_1} \ldots c_m^{n_m}
\delta_{n_1}(\sum r_{1t} x_t  - r_{10}) \ldots
\delta_{n_m}(\sum r_{mt} x_t  - r_{m0})
$$
where the sum is over $n_1 + \ldots + n_m = n$. This proves what we want.
\end{proof}

\end{document}
